\documentclass[11pt]{amsart}

\usepackage[margin=1in]{geometry}
\usepackage{amsmath,amssymb,amsthm,mathtools}
\usepackage{microtype}
\usepackage[T1]{fontenc}
\usepackage{lmodern}
\usepackage{hyperref}

\hypersetup{colorlinks=true,linkcolor=blue,citecolor=blue,urlcolor=blue}

\theoremstyle{plain}
\newtheorem{theorem}{Theorem}

\newcommand{\R}{\mathbb{R}}

\title{Wiener-Mean Laguerre Moments\\
of the Band-Limited Stationary Process}
\author{Stephen Crowley}
\date{\today}

\begin{document}

\maketitle

\begin{theorem}
Let $X$ be a real-valued stationary process on $\R$ with spectral density
\[
\widehat{X}_W(\xi) = \tfrac{1}{2}\mathbf{1}_{[-1,1]}(\xi)
\]
and Wiener covariance $K(h) = \sin(h)/h$. Define the generalized Laguerre functional
\[
L_n[X](u) = \tfrac{1}{2}\sum_{j=0}^{2n}(-1)^{j+n}\binom{2n}{j}\,X^{(j)}(u)\,X^{(2n-j)}(u),\qquad n \geq 0.
\]
Then
\[
\langle L_n[X]\rangle_W \;=\; \frac{4^n}{2(2n+1)}.
\]
\end{theorem}

\begin{proof}
By the reality of $X$, the spectral measure $\Phi$ satisfies $\Phi(-d\xi) = \overline{\Phi(d\xi)}$, so the Wiener-mean of $\Phi(d\xi)\Phi(d\eta)$ is supported on the anti-diagonal $\xi + \eta = 0$ with density $\mu(d\xi) = \tfrac{1}{2}\mathbf{1}_{[-1,1]}\,d\xi$. Substituting the Fourier representation $X^{(j)}(u) = \int(i\xi)^j e^{iu\xi}\,\Phi(d\xi)$ into the Laguerre definition, applying the binomial identity
\[
\sum_{j=0}^{2n}(-1)^j\binom{2n}{j}\xi^j\eta^{2n-j} = (\xi-\eta)^{2n}
\]
with $i^{2n} = (-1)^n$ canceling the alternating sign, and taking Wiener mean:
\[
\langle L_n[X]\rangle_W
= \tfrac{1}{2}\int_{-1}^{1}(2\xi)^{2n}\,\mu(d\xi)
= \tfrac{4^n}{2}\cdot\tfrac{1}{2}\int_{-1}^{1}\xi^{2n}\,d\xi
= \tfrac{4^n}{2}\cdot\tfrac{1}{2n+1}
= \frac{4^n}{2(2n+1)}.
\qedhere
\]
\end{proof}

\section*{Properties}

\begin{description}
  \item[(a) Stationarity.] The Wiener mean is independent of $u$.
  \item[(b) Automatic nonnegativity.] $(2\xi)^{2n} \geq 0$ pointwise and $\mu$ is positive, so the mean is $\geq 0$ for free.
  \item[(c) Universal fingerprint.] The values
  \[
  \left\{\frac{4^n}{2(2n+1)}\right\}_{n \geq 0} = \left\{\tfrac{1}{2},\,\tfrac{2}{3},\,\tfrac{8}{5},\,\tfrac{64}{21},\,\tfrac{128}{9},\,\ldots\right\}
  \]
  depend only on the spectral density being $\tfrac{1}{2}\mathbf{1}_{[-1,1]}$. Any band-limited stationary process with this density produces the same sequence.
  \item[(d) Mean-versus-path gap.] Sample-path positivity $L_n[X](u) \geq 0$ pointwise (which the Riemann Hypothesis requires) is strictly stronger than Wiener-mean positivity. The gap is the chaos correction; the proof of the Riemann Hypothesis closes sample-path positivity directly via the autocorrelation factorization $F_{n,\varepsilon} = g_{n,\varepsilon} \star g_{n,\varepsilon}$ and Bochner's theorem.
\end{description}

\end{document}
